\documentclass[a4paper,12pt]{article}

% Packages for formatting, math, and graphics
\usepackage{amsmath, amssymb, amsthm}
\usepackage{geometry}
\usepackage{booktabs}   % Professional tables
\usepackage{graphicx}   % Inclusion of images
\usepackage{pst-knot}   % For drawing knots (Requires XeLaTeX or latex->dvips->ps2pdf)
\usepackage{hyperref}   % For hyperlinks and cross-references
\usepackage{fancyhdr}   % Headers and footers
\usepackage{titling}    % Custom title

% Page setup
\geometry{margin=1in}
\hypersetup{
    colorlinks=true,
    linkcolor=blue,
    citecolor=red,
    urlcolor=cyan,
    pdftitle={Arithmetic Expression Geometry and Knot Theory},
    pdfauthor={Mingli Yuan}
}

% Theorem environments
\newtheorem{theorem}{Theorem}[section]
\newtheorem{definition}{Definition}[section]
\newtheorem{conjecture}{Conjecture}[section]
\newtheorem{observation}{Observation}[section]

% Title Information
\title{\textbf{Geometry of Arithmetic Expressions:}\\
\large Flows, Global Arithmetic Torsion, and the Topology of Knots}
\author{Mingli Yuan}
\date{\today}

\begin{document}

\maketitle

\begin{abstract}
This paper unifies the theory of Arithmetic Expression Geometry (AEG) with the topological study of knots. We introduce the concept of arithmetic expressions as geometric flows generated by orthogonal addition and multiplication fields. Within this framework, we define \textit{Global Arithmetic Torsion} ($\tau$), a measure of non-commutativity in arithmetic paths. We establish the \textbf{Triple Identity}, which links the algebraic definition of torsion to geometric area integrals within an Accumulative Commutative Space (ACS).

Applying this theory to knot groups, we demonstrate that mapping knot generators to arithmetic operators ($\otimes_t, \oplus_1$) transforms knot relators into closed loops in the expression space. We show that: (1) The closure condition of these loops recovers the Alexander polynomial $\Delta(t)=0$; and (2) The Global Arithmetic Torsion of these relator paths factors into $\Delta(t)$ and a cyclotomic component $(t^K-1)$, suggesting a deep connection between the non-commutativity of arithmetic operations and the topological symmetries of knot complements. Extensive computational results for knots $4_1, 6_2, 6_3,$ and $7_6$ are provided to support these findings.
\end{abstract}

\tableofcontents
\newpage

\section{Introduction: Arithmetic as Geometry}

The interplay between discrete algebra and continuous geometry has long fascinated mathematicians. Arithmetic Expression Geometry (AEG) proposes a novel perspective: arithmetic expressions are not merely static syntactic trees but dynamic paths in a geometric space.

\subsection{The Geometric Flow of Arithmetic}
In the AEG framework, elementary operations are interpreted as movements along orthogonal vector fields on a Riemann surface $\mathcal{M}$:
\begin{itemize}
    \item \textbf{Addition ($\oplus_\mu$):} Corresponds to a translation or "drift" flow, governed by parameter $\mu$.
    \item \textbf{Multiplication ($\otimes_\lambda$):} Corresponds to a scaling or "expansion" flow, governed by parameter $\lambda$.
\end{itemize}

We formalize this using a \textbf{Flow Equation}. For a scalar field assignment $a: \mathcal{M} \to \mathbb{R}$, the evolution of the value along a path parameterized by arc length $s$ and angle $\theta$ is given by:
\begin{equation}
    \frac{da}{ds} = \mu \cos \theta + a \lambda \sin \theta
\end{equation}
Or in its coordinate-free Eikonal form:
\begin{equation}
    ||\nabla a|| = \sqrt{\mu^2 + a^2 \lambda^2}
\end{equation}
This equation governs how values propagate through the \textit{First Kind Arithmetic Expression Space} $\mathfrak{E}_1$, which can be modeled on the upper half-plane with a hyperbolic metric.

\subsection{From Geometry to Topology}
This paper applies the AEG framework to Knot Theory. By treating the generators of a knot group as arithmetic operators acting on this space, knot relators (paths that must equal the identity) become closed loops. We investigate the geometric properties of these loops, specifically utilizing \textbf{Arithmetic Torsion} to measure the "gap" caused by the non-commutativity of operations.

\section{Global Arithmetic Torsion and the Triple Identity}

A central concept in AEG is that arithmetic operations do not commute ($a \times b + c \neq a + c \times b$). This non-commutativity generates a "torsion" in the geometric space.

\subsection{Algebraic Definition}
Let $\gamma$ be an arithmetic path (a sequence of operations). Let $\bar{\gamma}$ be the \textit{reversed-sequence path}, where the exact sequence of literal operations in $\gamma$ is applied in reverse order.
\begin{definition}[Global Arithmetic Torsion]
The global arithmetic torsion $\tau(\gamma)$ is defined as the difference between the evaluation of the path and its reversal:
\begin{equation}
    \tau(\gamma) = \nu(\gamma) - \nu(\bar{\gamma})
\end{equation}
where $\nu(\cdot)$ denotes the evaluation function.
\end{definition}

\subsection{The Accumulative Commutative Space (ACS)}
To quantify this torsion geometrically, we project the path onto an \textbf{Accumulative Commutative Space (ACS)}, a Euclidean plane where coordinates $(A, M)$ represent the accumulated additive ($\sum \mu$) and multiplicative ($\sum \lambda$) parameters. In this space, $\gamma$ and $\bar{\gamma}$ share start and end points, enclosing a region $\Sigma_\gamma$.

\subsection{The Triple Identity Theorem}
One of the foundational results of AEG is the Triple Identity, which bridges the algebraic and geometric views.

\begin{theorem}[The Triple Identity]
The following three quantities are equivalent for a closed structure in the ACS:
\begin{enumerate}
    \item \textbf{Algebraic Form:} The difference in evaluation, $\tau_{alg}(\gamma) = \nu(\gamma) - \nu(\bar{\gamma})$.
    \item \textbf{Interior Integral Form:} The weighted area integral over the region $\Sigma_\gamma$ enclosed by the paths in the ACS:
    \[ \tau_{int}(\gamma) = \iint_{\Sigma_\gamma} e^M \, dM \wedge dA \]
    \item \textbf{Boundary Integral Form:} The line integral along the boundary $\partial \Sigma_\gamma$:
    \[ \tau_{bound}(\gamma) = \oint_{\partial \Sigma_\gamma} e^M \, dA \]
\end{enumerate}
\end{theorem}

This theorem confirms that arithmetic torsion is not just a numerical artifact but a measurable geometric quantity (an $e^M$-weighted area) induced by the path's trajectory.

\section{Arithmetic Interpretation of Knot Groups}

We now apply this framework to the fundamental group of a knot complement. A knot group is typically presented as $\langle g_1, \dots, g_n \mid r_1, \dots, r_m \rangle$.

\subsection{The Mapping}
We map group generators to operators acting on an indeterminate variable $t$ (the parameter of the Alexander polynomial):
\begin{itemize}
    \item \textbf{Multiplicative Generator ($a$):} $a \mapsto \otimes_t$ (Multiply by $t$).
    \item \textbf{Additive Generator ($b$):} $b \mapsto \oplus_1$ (Add 1).
    \item \textbf{Inverses:} $A = a^{-1} \mapsto \oslash_t$ (Divide by $t$), $B = b^{-1} \mapsto \ominus_1$ (Subtract 1).
\end{itemize}

\subsection{Case Study: The Figure-Eight Knot ($4_1$)}

The knot $4_1$ has the presentation $\langle a, b \mid abbbaBAAB = 1 \rangle$.
The relator path is $\gamma_R = abbbaBAAB$.

\begin{figure}[h]
    \centering
    \begin{pspicture}(-2,-2)(2,2)
        \psKnot[linewidth=3pt,linecolor=blue](0,0){4-1}
    \end{pspicture}
    \caption{The Figure-Eight Knot ($4_1$). (Requires PSTricks compilation)}
    \label{fig:knot_4_1}
\end{figure}

\subsubsection{Derivation of the Alexander Polynomial}
Applying the arithmetic sequence corresponding to $abbbaBAAB$ to an initial value $x$:
\begin{align*}
    \nu(\gamma_R)(x) &= a(b(b(b(a(B(A(A(B(x))))))))) \\
    &= a(b(b(b(a(B(A(A(x-1)))))))) \\
    &= \dots \quad (\text{intermediate steps omitted}) \\
    &= x - (t^2 - 3t + 1)
\end{align*}
The closure condition $\nu(\gamma_R)(x) = x$ implies:
\begin{equation}
    t^2 - 3t + 1 = 0
\end{equation}
This is exactly the Alexander polynomial $\Delta_{4_1}(t) = 0$.

\subsubsection{Computational Verification}
This direct correspondence has been verified for numerous knots up to 11 crossings, including $6_2, 6_3, 7_6, 7_7, 8_2$, etc., suggesting that the AEG space naturally encodes the abelian invariants of the knot complement.

\section{Experimental Results: Torsion and Cyclotomic Factors}

We computed the Global Arithmetic Torsion $\tau(\gamma_R)(t) = p(t) - q(t)$ for various knots, where $p(t) = \nu(\gamma_R)(0)$ and $q(t) = \nu(\bar{\gamma}_R)(0)$.
The results reveal a consistent structure:
\begin{equation} \label{eq:torsion_factored_form}
\tau(\gamma_R)(t) = \sigma_{\text{eff}} \cdot \frac{\Delta(t)(t^K-1)}{t^k}
\end{equation}
This implies the torsion vanishes not only when $\Delta(t)=0$, but also when $t$ is a $K$-th root of unity. This links knot topology to cyclotomic fields.

\subsection{Summary Tables}

In the tables below, $\Phi_n(t)$ denotes the $n$-th cyclotomic polynomial (e.g., $\Phi_1 = t-1, \Phi_2 = t+1, \Phi_4 = t^2+1$).

\subsubsection{Knot $4_1$ ($\Delta(t) = t^2 - 3t + 1$)}
\begin{table}[htbp]
\centering
\caption{Torsion Analysis for Knot $4_1$}
\small
\begin{tabular}{p{3.5cm} p{1.8cm} p{1.8cm} p{3.5cm} p{3cm}}
\toprule
\textbf{Relator} & \textbf{$p(t)$} & \textbf{$q(t)$} & \textbf{Torsion $\tau(t)$} & \textbf{Factors} \\
\midrule
\texttt{abbbaBAAB} & $-\Delta(t)$ & $\frac{-\Delta(t)}{t^2}$ & $\frac{-\Delta(t)(t^2-1)}{t^2}$ & $\Phi_1 \Phi_2$ \\
\texttt{bbbaBAABa} & $\frac{-\Delta(t)}{t}$ & $\frac{-\Delta(t)}{t}$ & $0$ & None \\
\bottomrule
\end{tabular}
\end{table}

\subsubsection{Knot $6_2$ ($\Delta(t) = t^4 - 3t^3 + 3t^2 - 3t + 1$)}
\begin{table}[htbp]
\centering
\caption{Torsion Analysis for Knot $6_2$}
\small
\begin{tabular}{p{5cm} p{2cm} p{4cm} p{3cm}}
\toprule
\textbf{Relator Type} & \textbf{$p(t)$} & \textbf{Torsion $\tau(t)$} & \textbf{Cyclotomic} \\
\midrule
\texttt{aaBAbbbAAbbbABaaBBB} & $\frac{-\Delta}{t^2}$ & $0$ & None \\
\texttt{bbbAAbbbABaaBBBaaBA} & $\frac{-\Delta}{t^3}$ & $\frac{\Delta(t)(t^2-1)}{t^3}$ & $\Phi_1 \Phi_2$ \\
\texttt{aabbbaBAABBBAABabbb} & $\frac{-\Delta}{t}$ & $\frac{-\Delta(t)(t^2-1)}{t^3}$ & $\Phi_1 \Phi_2$ \\
\texttt{abbbaabbbaBAABBBAAB} & $-\Delta$ & $\frac{-\Delta(t)(t^4-1)}{t^4}$ & $\Phi_1 \Phi_2 \Phi_4$ \\
\bottomrule
\end{tabular}
\end{table}

\subsubsection{Knot $6_3$ ($\Delta(t) = t^4 - 3t^3 + 5t^2 - 3t + 1$)}
Knot $6_3$ is amphicheiral, and this symmetry is reflected in the variety of torsion results.

\begin{table}[htbp]
\centering
\caption{Torsion Analysis for Knot $6_3$}
\small
\begin{tabular}{p{5cm} p{2cm} p{4cm} p{3cm}}
\toprule
\textbf{Relator Type} & \textbf{$p(t)$} & \textbf{Torsion $\tau(t)$} & \textbf{Cyclotomic} \\
\midrule
\texttt{baabbaBAABabbaabABB...} & $\Delta(t)$ & $\frac{\Delta(t)(t^4-1)}{t^4}$ & $\Phi_1 \Phi_2 \Phi_4$ \\
\texttt{aBAABBAbaabbbaabABB...} & $\frac{\Delta}{t^2}$ & $0$ & None \\
\texttt{aaBAbbAAbaBBabaBB...} & $\frac{\Delta}{t}$ & $\frac{\Delta(t)(t^2-1)}{t^3}$ & $\Phi_1 \Phi_2$ \\
\texttt{bABBAbABBAbaabbaB...} & $\frac{\Delta}{t^4}$ & $\frac{-\Delta(t)(t^4-1)}{t^4}$ & $\Phi_1 \Phi_2 \Phi_4$ \\
\bottomrule
\end{tabular}
\end{table}

\section{Discussion: The Rationality Hypothesis}

The emergence of the factor $(t^K-1)$ in the torsion term is physically significant. It suggests that the "forward" flow $\gamma$ and the "reverse" flow $\bar{\gamma}$ differ by a geometric phase that is rational or commensurate with the periodicity of the underlying covering space.

\subsection{Periodicity and Fundamental Domains}
We hypothesize that the relator path $\gamma_R$ traces the boundary of a fundamental domain in the knot complement. When lifted to the infinite cyclic cover, the difference between the forward and reverse arithmetic evaluations ($k_p$ vs $k_q$) reflects the relative "winding" or "phase depth" of these paths in the $t$-graded structure. The integer difference $K = |k_q - k_p|$ implies a "simple ratio" commensurability, essential for the discrete group $\Gamma$ to act consistently on the universal cover.

\subsection{Comparison with Random Paths}
Computational experiments comparing relator paths against randomly shuffled paths (preserving operation counts) show that:
\begin{enumerate}
    \item Random paths do not generate $\Delta(t)$ in their evaluation.
    \item The difference $p_{val}(t) - p_{shuffled}(t)$ does not factor neatly into cyclotomic polynomials.
\end{enumerate}
This confirms that the algebraic structure observed (Triple Identity satisfaction and Cyclotomic factorization) is an intrinsic property of the topological knot relator, not a statistical artifact of the operation count.

\section{Conclusion}
By viewing arithmetic expressions as geometric flows, we have established a framework where knot relators naturally generate the Alexander polynomial. The concept of Global Arithmetic Torsion, grounded in the Triple Identity, reveals that the non-commutativity of the knot group is structurally linked to cyclotomic fields and roots of unity. This suggests that Arithmetic Expression Geometry provides a faithful geometric representation of the homological symmetries inherent in knot topology.

\begin{thebibliography}{9}
\bibitem{aeg} Yuan, M. \textit{Geometry of Arithmetic Expressions: I. Basic Concepts and Unsolved Problems}. (Draft), 2025.
\bibitem{alexander} Alexander, J. W. \textit{Topological invariants of knots and links}. Trans. Amer. Math. Soc. 30 (1928), 275–306.
\bibitem{fox} Fox, R. H. \textit{Free differential calculus}. Ann. of Math. (2) 57 (1953), 547–560.
\bibitem{repo} Yuan, M. \textit{Knot Alexander Repository}. \url{https://github.com/mountain/knot-alexander}.
\end{thebibliography}

\end{document}